\documentclass{article}
\usepackage{amsmath}
\usepackage{amsfonts}
\usepackage{graphicx}

\newcommand{\transf}[1]{\mathcal H(#1)}
\newcommand{\freq}[1]{H(#1)}
\newcommand{\magres}[1]{|\freq{#1}|}

\begin{document}
\title{Mapping Filters from Analog to Digital with Tustin's Method}
\author{Edgar Jose Donoso Mansilla}
\date{}
\maketitle

% TODO: Review the paper and ask about the topic to the prof

\section{Motivation}
Real world systems are often continuous, therefore, there is a lot of pre-existing
literature and techniques to handle continuous signals via continuous means
(e.g. analog circuit based filters). As computing and integrated chips became
a more widely used tool, there was a need to convert some well studied filters
and systems to discrete time despite there being DSP native techniques. The key
goal of these conversion methods is to map values from the $s$-plane that is obtained
via Laplace transforms to the $z$-plane that is obtained from the $z$-transform.

\section{Tustin's Method Overview}
Tustin's Method, also known as the Bilinear Transform, is a method for converting
analog filters to digital filters. The formal derivation of this method is as
follows:

Let $s \in \mathbb C$ and $T$ be the integration step (to be discussed later).
$$
	z = e^{sT} = \frac{e^{sT/2}}{e^{-sT/2}}
$$

From here, we can approximate the exponential using the Taylor series of $e^x$:
$$
	e^{sT/2} = 1 + \frac{sT/2}{1!} + \frac{(sT/2)^2}{2!} + \frac{(sT/2)^3}{3!} + \dots
$$

By taking the first 2 terms we can obtain:
$$
	z = \frac{1 + sT/2}{1 - sT/2} \rightarrow s = \frac{2(z - 1)}{T(z + 1)}
$$

\subsection{Understanding the meaning of $T$}
While this shows the algebraic manipulation that explains the mapping, it is
better understood by looking at integration that this approximation is doing:

Let a continuous system:
$$
	H(s) = 1/s
$$

Taking the inverse Laplace transform of that system:
$$
	y'(t) = x(t)
$$

Integrate to solve for $y(t)$:
$$
	y(t) = \int_0^t{x(t) \, dt}
$$

Then redefine the system as steps based on $T$:
$$
	y(kT) = \int_{kT-T}^{kT}{x(t) \, dt} + y(kT - T)
$$

Recall Trapezoidal Integration:
$$
	\int_a^b{f(x) \, dx} \approx (b - a)\frac{f(a) + f(b)}{2}
$$

And thus the following discrete system is obtained:
$$
	y_k = \frac{T}{2}(x_{kT} + x_{kT - T}) + y_{kT - T}
$$

After getting the $z$-transform of this system\footnote{Or in the context of
	lecture, obtain the discrete transfer function $\mathcal H(z)$}:
$$
	\transf{z} = \frac{T (z + 1)}{2 (z - 1)}
$$

We also know that this is supposed to be the same as the transfer function
of our continuous system, therefore:
$$
	\frac{1}{s} \approx \frac{T (z + 1)}{2 (z - 1)} \rightarrow z \approx \frac{1 + sT/2}{1 - sT/2}
$$

\section{Driving an oscillation through digital input}
A common classical control theory problem involves manipulating the frequency of
an oscillating mass. However, sometimes the controller has to be implemented through
an IC or microcontroller. For example, this system with input force $x(t)$:
$$
	y''(t) = x(t) \rightarrow H(s) = \frac{1}{s^2}
$$
It doesn't oscillate at all, therefore the input force can be used as a controller.

Then we can use Tustin's Method to discretize our system:
$$
	\transf{z} = \frac{4(z - 1)^2}{T^2(z + 1)^2}
$$

Notice how the system went from having a pole at $s = 0$ to haivng poles at
$z = \pm 1$. This is one of the side-effects of using Tustin's method.

After rearranging to get the digital system\footnote{This technique could be used
	to obtain synthesized instruments through physical principles. Albeit it gets
	complicated past the simple oscillation case}:
$$
	y_t = \frac{T^2}{4}(x_t + 2x_{t-1} + x_{t-2}) + 2y_{t-1} - y_{t-2}
$$

In the ideal case an arbitrary force can be applied to the controller at each
time-step, a simple control law to get the desired frequency $r$ can be obtained:
$$
	x_t = C_p(r - y_{t-1}) - C_d(\frac{y_{t-1} + y_{t-2}}{2})
$$

In this controller we are using $C_p$ to approach our desired frequency while
$C_d$ will dampen the system. Further analysis would be needed to prove the
system remains stable under various coefficient values, however, that is
beyond the scope of this short paper.

\end{document}
